\section{Introduction}
\label{sec:introduction}

Avocado (\emph{Persea americana}) is a high-value crop with growing importance in the global agricultural economy, yet its production is constantly threatened by leaf diseases and pest damage, including leaf scorch, leaf spot, algal rust, powdery mildew, and sap-feeding mealybugs, which can severely reduce fruit yield and quality. Early and accurate diagnosis is therefore essential for effective crop management, but in practice it still relies on visual inspection by human experts, a process that is time-consuming, subjective, and impossible to scale to large orchards.

Deep learning, particularly convolutional neural networks (CNNs), has emerged as a promising alternative, often reporting remarkable accuracy on controlled benchmarks such as PlantVillage. Translating that promise to real farms, however, exposes four gaps that controlled benchmarks tend to hide. First, most high-accuracy models are computationally heavy, requiring GPUs that are impractical for edge devices such as smartphones or a Raspberry Pi in the field. Second, collecting and annotating thousands of images per disease class is expensive and labour-intensive, while deep models conventionally assume exactly that scale of labelled data is available. Third, a model that reports a class label without any notion of its own uncertainty is unsafe to act on in high-stakes decisions, and offers no signal for when a diagnosis should be deferred to a human expert. Fourth, a black-box prediction, however accurate, does little to build the trust of end-users such as farmers and agronomists, who need to see \emph{why} a leaf was flagged as diseased before they act on it.

This paper addresses all four gaps around a single proposed model, which we name \textbf{BoVien}\footnote{From Vietnamese \emph{bơ} (avocado) + \emph{viên} (a small unit / device).}, with four contributions.

\begin{enumerate}
    \item \textbf{AvocadoLeaf dataset.} A new dataset of 1,861 images covering five disease categories and a healthy class, captured first-hand under diverse real-world conditions (variable lighting, multiple angles, natural backgrounds) rather than in a controlled studio setting. The dataset exhibits moderate class imbalance (max/min ratio 3.96) and fine-grained inter-class similarity (Silhouette score 0.0067, intra/inter distance ratio 1.59), reflecting the genuine difficulty of field-based diagnosis rather than an artificially easy benchmark.

    \item \textbf{Model BoVien, a lightweight architecture for edge deployment.} Selected from a systematic benchmark of 15 architectures spanning CNN, hybrid CNN-Transformer, and pure-transformer families (Section~\ref{sec:experiments}), BoVien (a ConvNeXt-V2 backbone with selective layer fine-tuning) reaches the best accuracy-size-latency trade-off, with 92.14\% test accuracy and 0.900 macro F1 at only 3.4M parameters and 2.30\,ms per image on a single GPU.

    \item \textbf{Few-shot learning as a data-efficient fallback.} To probe how the approach degrades when labelled data is scarce, we compare standard fine-tuning against Prototypical Networks (ProtoNet) and Relation Networks. ProtoNet more than doubles fine-tuning's accuracy at every shot level tested, reaching 57.4\% accuracy from only 20 labelled images per class (versus 28.6\% for fine-tuning), showing that metric-based few-shot learning is a viable fallback when a new class or a new deployment site starts with little labelled data.

    \item \textbf{Uncertainty estimation and out-of-distribution (OOD) detection.} Using Monte Carlo Dropout, BoVien separates in-distribution inputs from out-of-distribution ones with high reliability, reaching AUC 0.994 on far-OOD (unrelated internet images) and AUC 0.985 on the harder near-OOD case (leaves of other crop species), giving a practical signal for when a prediction should be flagged for human review. Its raw softmax confidence, by contrast, is only moderately calibrated (ECE = 0.0625), which we discuss as a limitation in Section~\ref{sec:discussion} rather than overstate here. This OOD test asks whether BoVien recognises \emph{off-task} inputs (non-leaf images, leaves of other crops) as unfamiliar; it is a distinct question from whether BoVien still recognises \emph{avocado leaves collected at a new site}, which Section~\ref{subsec:real_cross_domain} answers separately and less favourably.
\end{enumerate}

Figure~\ref{fig:pipeline_overview} summarises how these four contributions fit together, from field data collection through architecture selection to the explainability, robustness, and edge-deployment evaluation that follows.

\begin{figure*}[t]
\centering
\resizebox{\textwidth}{!}{%
\begin{tikzpicture}[
    font=\sffamily,
    every node/.style={align=center},
    phase/.style={rounded corners=4pt, minimum height=2.6cm, text width=2.05cm,
                  text=white, font=\sffamily\bfseries\small, drop shadow},
    bullet/.style={text=black!75, font=\sffamily\scriptsize, align=left, text width=2.15cm},
    arr/.style={-{Stealth[length=2.8mm]}, line width=1pt, draw=black!55},
]

% ============================================================
% LEFT: Dataset Construction
% ============================================================
\node[draw=black!45, fill=blue!4, rounded corners=2pt, minimum width=3.3cm, minimum height=1.7cm]
    (collect) at (0,0) {\faCamera\ Field collection\\\scriptsize 3 provinces, Central Highlands};

\node[draw=black!45, fill=green!4, rounded corners=2pt, minimum width=3.3cm, minimum height=1.5cm,
      below=0.32cm of collect] (annotate) {\faUsers\ Expert annotation\\\scriptsize Cohen's $\kappa$ agreement};

\node[draw=black!45, fill=orange!5, rounded corners=2pt, minimum width=3.3cm, minimum height=1.5cm,
      below=0.32cm of annotate] (stats) {\faDatabase\ AvocadoLeaf dataset\\\scriptsize 1{,}861 imgs, 6 classes};

\draw[arr] (collect) -- (annotate);
\draw[arr] (annotate) -- (stats);

\node[font=\sffamily\bfseries\normalsize, above=0.55cm of collect] (leftlabel) {Dataset Construction};

\node[draw=black!60, dashed, rounded corners=3pt, fit=(leftlabel)(collect)(annotate)(stats),
      inner xsep=0.32cm, inner ysep=0.3cm] (dataleft) {};

% ============================================================
% RIGHT: Systematic Study
% ============================================================
\node[phase, top color=blue!55, bottom color=blue!75!black,
      right=2.9cm of dataleft, yshift=0.55cm] (bench) {\faThLarge\\[2pt]Benchmark\\15 Archi-\\tectures};

\node[phase, top color=violet!60, bottom color=violet!85!black,
      right=0.42cm of bench, minimum height=2.9cm, line width=1.4pt, draw=black] (bovien) {\faStar\\[2pt]Model\\BoVien\\(Selected)};

\node[phase, top color=teal!55, bottom color=teal!80!black,
      right=0.42cm of bovien] (xai) {\faSearch\\[2pt]Explain-\\ability};
\node[bullet, below=0.2cm of xai] (xaib) {$\bullet$ LIME\\$\bullet$ VLM caption\\$\bullet$ Faithfulness};

\node[phase, top color=green!45!black!15!green, bottom color=green!55!black,
      right=0.42cm of xai] (rob) {\faUserShield\\[2pt]Robust-\\ness};
\node[bullet, below=0.2cm of rob] (robb) {$\bullet$ Uncertainty\\$\bullet$ Few-shot\\$\bullet$ Cross-dataset};

\node[phase, top color=orange!65, bottom color=orange!85!black,
      right=0.42cm of rob] (deploy) {\faMicrochip\\[2pt]Edge De-\\ployment};
\node[bullet, below=0.2cm of deploy] (deployb) {$\bullet$ Real-time\\$\bullet$ Latency\\$\bullet$ On-device};

\draw[arr] (bench) -- (bovien);
\draw[arr] (bovien) -- (xai);
\draw[arr] (xai) -- (rob);
\draw[arr] (rob) -- (deploy);

\node[font=\sffamily\bfseries\normalsize] (rightlabel) at ($(bench.north)!0.5!(deploy.north)+(0,1.0)$) {Systematic Study};

\node[draw=black!60, rounded corners=3pt,
      fit=(rightlabel)(bench)(bovien)(xai)(rob)(deploy)(xaib)(robb)(deployb),
      inner xsep=0.4cm, inner ysep=0.35cm] (studyright) {};

% ============================================================
% Data-preprocessing connector, centred in the gap between panels
% ============================================================
\node[font=\sffamily\bfseries\scriptsize, text width=1.4cm, align=center, rotate=90]
    (prep) at ($(dataleft.east)!0.5!(studyright.west)$) {Data Preprocessing};
\draw[arr, line width=1.3pt]
    ($(dataleft.east)!0.5!(studyright.west)+(0,1.35cm)$) --
    ($(dataleft.east)!0.5!(studyright.west)+(0,-1.35cm)$);

% ============================================================
% Outer frame + title
% ============================================================
\node[font=\sffamily\bfseries\Large] (title) at ($(dataleft.north west)!0.5!(studyright.north east)+(0,0.55cm)$)
    {Overall Study Pipeline of Model BoVien};

\node[draw=black!70, line width=1pt, rounded corners=2pt,
      fit=(title)(dataleft)(studyright), inner xsep=0.4cm, inner ysep=0.3cm] (outer) {};

\draw[black!50] (outer.north west |- title.south) ++(0.15cm,-0.15cm) -- (outer.north east |- title.south) ++(-0.15cm,-0.15cm);

\end{tikzpicture}%
}
\caption{Overview of the overall study pipeline. On the left, AvocadoLeaf dataset construction (field collection across three Central Highlands provinces, expert annotation with inter-annotator agreement, yielding 1{,}861 quality-filtered images across 6 classes). On the right, the systematic study built on top of it, benchmarking 15 architectures, selecting the proposed Model BoVien, then evaluating it for explainability, robustness (uncertainty/OOD, few-shot, cross-dataset), and edge-deployment readiness.}
\label{fig:pipeline_overview}
\end{figure*}

The remainder of this paper is organised as follows. Section~\ref{sec:related_work} reviews related work in plant disease classification, few-shot learning, and uncertainty estimation. Section~\ref{sec:dataset_construction} describes the AvocadoLeaf dataset construction process, and Section~\ref{sec:dataset_analysis} analyses its statistics, photometric properties, and intrinsic difficulty. Section~\ref{sec:proposed_model} details the 15-architecture benchmark and the resulting BoVien architecture. Section~\ref{sec:experiments} reports the full experimental results, covering architecture comparison, explainability, few-shot learning, uncertainty estimation, cross-dataset evaluation, and real-time inference. Section~\ref{sec:discussion} discusses findings and limitations, and Section~\ref{sec:conclusion} concludes the paper and outlines future work.